%%%%%%%%%%%%%%%%%%%%%%%%%%%%%%%%%%%%%%%%%%%%%%%%%
\section{Power counting}
\label{sec:powercounting}

\subsection{Superficial UV divergences}

In order to identify all UV divergences from any Feynman diagram contributing to the quasi-PDFs, we start with a set of general diagrams constructed from the lower order Feynman diagrams by adding one more gluon to them. We will then examine how this additional gluon could change the UV divergence of the original diagrams. Specifically, we introduce and study the change of divergence index $\Delta\omega_3$ and $\Delta\omega_4$ corresponding to the 3-D integration and 4-D integration of corresponding loop, respectively. A sufficient condition for quasi-PDFs to be renormalizable is that $\Delta\omega_3\le 0$ and $\Delta\omega_4\le 0$ are satisfied for all cases, but it is not necessary as we will argue. We divide our discussions into five distinctive cases:

%%%%%%%%%%%%%%%%%%%%%%%%%%%%%%%
\begin{figure}[t]
\begin{center}
\includegraphics[width=3in]{addonegluon}
\caption{\label{fig:addonegluon}
Some loop diagrams, with an additional gluon (denoted as dotted curves) attachments, that could contribute to the quasi-PDFs at higher loop orders.}
\end{center}
\end{figure}
%%%%%%%%%%%%%%%%%%%%%%%%%%%%%%%

\textbf{Case I:} Only one end of the gluon is attached to parton lines of a loop, like Fig.~\ref{fig:addonegluon}(a). In this case we have one more propagator and one more QCD vertex in the loop, which results in $\Delta\omega_3=-1$ and $\Delta\omega_4=-1$. Thus a linear divergence can be changed to a logarithmic divergence, and a logarithmic divergent diagram is changed to be finite.

\textbf{Case II:} Only one end of the gluon is attached to gauge link of a loop, like Fig.~\ref{fig:addonegluon}(b). In this case, we have one more $z-$direction integration in coordinate space, which does not change the degree of divergence of 3-D integration but reduce the degree of divergence of 4-D integration by 1, such that $\Delta\omega_3=0$ and $\Delta\omega_4=-1$.

\textbf{Case III:} Both ends of the gluon are attached to parton lines of a loop, like Fig.~\ref{fig:addonegluon}(c). In this case we have three more propagators, two more QCD vertexes, and four more momentum integrations for the loop, which results in $\Delta\omega_3=-1$ and $\Delta\omega_4=0$.

\textbf{Case IV:} One end of the gluon is attached to parton lines of a loop, and the other end of the gluon is attached to gauge link of the loop, like Fig.~\ref{fig:addonegluon}(d). In this case we have two more propagators, one more QCD vertex, one more $z-$direction integration in coordinate space, and four more momentum integrations for the loop, which results in $\Delta\omega_3=0$ and $\Delta\omega_4=0$.

\textbf{Case V:} Both ends of the gluon are attached to gauge link of a loop, like Figs.~\ref{fig:addonegluon}(e,f). In this case we have one more propagator, two more $z-$direction integration in coordinate space, and four more momentum integrations for the loop. This results in $\Delta\omega_3= 1$ and $\Delta\omega_4=0$.

A few comments for the above power counting rules are in order.
1) We only concentrated on overall divergences but not on subdivergences, because subdivergences can be taken care by forest subtraction in the step of renormalization. 2) The change of divergence indexes discussed here are only for superficial divergence. It means that even if the power counting indicates that a diagram is divergent, it may not be really divergent due to other considerations. But, on the other hand, if the power counting indicates that a diagram is finite, it must be UV finite.

The \textbf{Case V} is dangerous for renormalization, because $\Delta\omega_3>0$ means that the number of potentially UV divergent topologies is not finite, and thus a finite number of renormalization constants may not be enough to remove all UV divergences. However, as we pointed out, the above power counting rules are only for superficial divergence. We will show in the next subsection that all Feynman diagrams of quasi-PDFs do not have overall UV divergence if only the 3-D integration is considered for any of its loop momenta. As a result, only $\Delta\omega_4$ is relevant for identifying real UV divergent topologies.

%%%%%%%%%%%%%%%%%%%%%%%%%%%%%%%
\subsection{Finiteness of the 3-D integration}

To demonstrate that the 3-D integration of Feynman diagrams contributing to the quasi-PDFs cannot generate a real UV divergence, we consider the ``gauge-link-irreducible" (GLI) diagrams.  Similar to the concept of one-particle-irreducible (1PI) diagrams, a GLI diagram is a diagram that is still connected even if all gauge links are cut (or removed).  For example, the diagrams (a), (b) and (c) in Fig.~\ref{fig:oneloop} are not the GLI diagrams, while the diagram (d) is.    Similarly, the diagrams (a), (c) and (d) in Fig.~\ref{fig:addonegluon} are the GLI diagrams, while the diagrams (b), (e) and (f) are not.   For studying the renormalization of quasi-PDFs, since the renormalizaiton of QCD Lagrangian is well known, we only need to study the UV properties of the general GLI diagrams as shown in Fig.~\ref{fig:connected}, where the dashed lines connecting to gauge link can be either gluon or quark (if it is at the end point of the gauge link).

%%%%%%%%%%%%%%%%%%%%%%%%%%%%%%%
\begin{figure}[h]
\begin{center}
\includegraphics[width=1in]{strongconnected}
\caption{\label{fig:connected}
A general ``gauge-link-irreducible" (GLI) diagram, where $l_0=q-l_1-\cdots-l_n$.}
\end{center}
\end{figure}
%%%%%%%%%%%%%%%%%%%%%%%%%%%%%%%
In Fig.~\ref{fig:connected}, we assume that the momentum flows into the blob is $q$,  $l_1,\cdots,l_n$ are $n$ loop momenta, and $l_0=q-l_1-\cdots-l_n$ is determined by the momentum conservation. Note that this assignment of momentum is always possible for a GLI diagram. Then,  the general diagram in Fig.~\ref{fig:connected} can be expressed as
\begin{align}\label{eq:connected}
e^{i q_{z} r_0} \prod_{j=1}^n \int_{r_{j-1}+a}^{r_{max}-a}d r_j \int \frac{d^4l_j}{(2\pi)^4} e^{i l_{jz} (r_j-r_0)} \mathcal{M}(q,l_1,\cdots,l_n),
\end{align}
where $\mathcal{M}(q,l_1,\cdots,l_n)$ denotes the blob combined with propagators of the $n+1$ lines connecting to the gauge link.  Since the GLI diagrams can be constructed from one-loop diagrams in Figs.~\ref{fig:oneloop} or~\ref{fig:oneloopg}  combined with insertions in \textbf{Cases I, III} and \textbf{IV} in the last subsection, their overall superficial UV divergence index $\omega\leq1$, no matter we apply the 3-D integration or 4-D integration to each loop momentum.

Now let us study the case in which only the 3-D integration is performed for $l_j$, while carrying out either 3-D or 4-D integration for other loop momenta.  The $\mathcal{M}(q,l_1,\cdots,l_n)$  has two kinds of dependence on $l_j$. One is the polynomial dependence in the numerator, for which we can simply decompose $l_j$ to $\bar{l}_{j}$ and $l_{jz}$. The other is in the propagator like $1/(l_j+k)^2$ where $k$ can be vanishing or depending on other loop momenta. We can decompose the propagator as
\begin{align}\label{eq:decomZ}
\begin{split}
\frac{1}{(l_j+k)^2}=& \frac{1}{\Delta-2k_z l_{jz}-l_{jz}^2}\\
=& \frac{1}{\Delta}+ \frac{2k_z l_{jz}}{\Delta^2}+\frac{(\Delta+4k_z^2+2k_zl_{jz}) l_{jz}^2}{(\Delta-2k_z l_{jz}-l_{jz}^2)\Delta^2},
\end{split}
\end{align}
where $\Delta=(\bar{l}_j+\bar{k})^2-k_z^2$ independent of $l_{jz}$. Based on dimensional counting, each additional $l_{jz}$ in the numerator will suppress the divergence index $\omega$ of the 3-D integration by one unit. Thus the last term in Eq.~\eqref{eq:decomZ} can be safely ignored as we are considering UV divergence from the 3-D integration of $l_j$.  Since the other terms factorize out the $l_{jz}$ dependence from $\mathcal{M}(q,l_1,\cdots,l_n)$, potential UV divergences of these terms are proportional to
\begin{align}
\int d l_{jz} e^{i l_{jz} (r_j-r_0)}  l_z^m \propto \delta^{(m)}(r_j-r_0),
\label{eq:deltam}
\end{align}
with $m$ being non-negative integer. As $r_j$ cannot equal to $r_0$ as defined in Eq.~\eqref{eq:connected}, $\delta^{(m)}(r_j-r_0)$ in Eq.~(\ref{eq:deltam}) vanishes before we take $a\to 0$ limit. Therefore, UV divergences of Fig.~\ref{fig:connected}, obtained by integrating out $\bar{l}_j$ and other loop momenta but fixing $l_{jz}$, eventually vanish after the integration of $l_{jz}$.

In summary, the overall UV divergences of GLI diagrams only come from the region where all loop momenta are large. Furthermore, the third term in Eq.~\eqref{eq:decomZ} is responsible for real UV divergences from the 4-D integration of the GLI diagrams.



%%%%%%%%%%%%%%%%%%%%%%%%%%%%%%%
\begin{figure}[h]
\begin{center}
\includegraphics[width=1.3in]{disconnected}
\caption{\label{fig:disconnect}
A diagram made of two ``gauge-link-irreducible'' (GLI) sub-diagrams.}
\end{center}
\end{figure}
%%%%%%%%%%%%%%%%%%%%%%%%%%%%%%%

Now let us study a diagram made of two GLI sub-diagrams as shown in Fig.~\ref{fig:disconnect}. This non-GLI diagram can be generated either from the diagram in Fig.~\ref{fig:twoloop} or by the insertion of \textbf{Cases II} and \textbf{ V} in the last subsection. Based on the power counting rules, this diagram has overall superficial UV divergence index $\omega\leq2$. When extracting overall UV divergences, we can follow the above discussion for each GLI sub-diagram, and found that the overall UV divergence of the combined diagram vanishes if we fix the $z$-component of any loop momentum. This conclusion can be easily generalized to any diagram made of more GLI sub-diagrams.


%%%%%%%%%%%%%%%%%%%%%%%%%%%%%%%
\begin{figure}[h]
\begin{center}
\includegraphics[width=1.1in]{twoloopbase}
\caption{\label{fig:twoloop}
A two-loop diagram that might give UV divergent contribution to quasi-quark PDF based on the power counting rules and building blocks in Figs.~\ref{fig:oneloop} and \ref{fig:oneloopg}.}
\end{center}
\end{figure}
%%%%%%%%%%%%%%%%%%%%%%%%%%%%%%%

We thus conclude that overall UV divergences of any diagram that contributes to quasi-PDFs can only come from the region where all loop momenta are large. An immediate consequence of this finding is that to identify UV divergent topologies, we only need to consider the value of $\Delta\omega_4$ when we are evaluating the five case insertions discussed in the last subsection.  Since we found that $\Delta\omega_4\le 0$ for all cases, there should be only finite number of topologies of UV divergent Feynman diagrams that could contribute to the quasi-PDFs, which we will present in the next subsection.

Let us conclude this subsection by explaining the key difference in UV behavior between quasi-PDFs and PDFs, which have UV divergence from the 3-D integration. As pointed out at the end of Sec.~\ref{sec:oneloop}, UV divergences of PDFs are obtained from a different 3-D integration of loop momentum, say $l_-$ and $\vec{l}_\perp$. We can certainly do a similar decomposition of propagators as that in Eq.~\eqref{eq:decomZ},
\begin{align}\label{eq:decomPlus}
\begin{split}
\frac{1}{(l+k)^2}=& \frac{1}{\hat{\Delta}+2l_+(l_-+k_-)}\\
=& \frac{1}{\hat{\Delta}}- \frac{2(l_-+k_-)l_+}{\hat{\Delta}^2}+\frac{4(l_-+k_-)^2l_+^2}{(\hat{\Delta}+2l_+(l_-+k_-))\hat{\Delta}^2},
\end{split}
\end{align}
where $\hat{\Delta}=2k_+(l_-+k_-)-(\vec{l}_\perp+\vec{k}_\perp)^2$, which is independent of $l_+$. We can also argue that, because $l_+$ is factorized out, the first two terms do not contribute after the integration of $l_+$. However, the last term in Eq.~\eqref{eq:decomPlus} can still generate the UV divergence from the  3-D integration. This is because the UV divergence for PDFs comes from the region where any loop momentum $l$ behaves as $(l_{+}, l_{-}, \vec{l}_{\perp})\sim(1,\lambda^2,\lambda)$ as $\lambda\to\infty$, and thus $l_- l_+ \sim l_\perp^2 \sim \lambda^2$. The boost invariance ensures that each $l_+$ in the numerator will be either accompanied with a ``$-$" momentum component in the numerator or a ``+'' momentum component in the denominator, and neither of these cases suppresses UV divergence index of the loop momentum integration.


%%%%%%%%%%%%%%%%%%%%%%%%%%%%%%%
\begin{figure}[t]
\begin{center}
\includegraphics[width=3in]{multiloopdiv}
\caption{\label{fig:multiloopdiv}
Three topologies of diagrams that could give UV divergent contributions to the quasi-quark PDFs.}
\end{center}
\end{figure}
%%%%%%%%%%%%%%%%%%%%%%%%%%%%%%%

%%%%%%%%%%%%%%%%%%%%%%%%%%%%%%%
\subsection{Divergent diagrams for quasi-quark PDFs}

Based on the above strategy and power counting rules, we can find out all UV divergent Feynman diagrams begin with a complete set of building blocks. To identify all UV divergent Feynman diagrams for quasi-quark PDFs, we need one more diagram in Fig.~\ref{fig:twoloop} to form the set of building blocks in addition to the one-loop diagrams in Figs.~\ref{fig:oneloop} and~\ref{fig:oneloopg}, because the superficial UV divergence of this two-loop diagram does not agree with that obtained from Fig.~\ref{fig:oneloop}(b) by inserting one more gluon like in \textbf{Case IV}.  Because we only need to consider 4-D integrations, non-vanishing $\xi_z$ in Fig.~\ref{fig:twoloop} ensures that this diagram is UV finite.

Using the above building blocks, we can generate all possible higher order Feynman diagrams. Among them, there are only three types of topologies that could give UV divergent contribution to the quasi-quark PDFs as shown in Fig.~\ref{fig:multiloopdiv}, in addition to those from the renormalizaiton of QCD Lagrangian. From the above discussion, we know that all of these three topologies can be UV divergent only in the region where all loop momenta go to infinity, thus UV divergent contributions are localized in coordinate space. An immediate consequence is that to make the diagrams in Fig.~\ref{fig:multiloopdiv} divergent, $r_2-r_1$ must go to 0. This finding also explains the result of the two-loop calculation in Ref.~\cite{Ji:2015jwa}, where one finds that UV divergence of quasi-quark PDF under DR is proportional to $\delta(1-z)$ in momentum space.

In addition to these three UV divergent topologies in Fig.~\ref{fig:multiloopdiv},  we also show some topologies that are UV finite in Fig.~\ref{fig:multiloopfin}.  Especially, the last diagram in Fig.~\ref{fig:multiloopfin} indicates that quasi-quark PDFs do not mix with quasi-gluon PDF under the renormalization to all orders in QCD perturbation theory.


%%%%%%%%%%%%%%%%%%%%%%%%%%%%%%%
\begin{figure}[h]
\begin{center}
\includegraphics[width=3in]{multiloopfin}
\caption{\label{fig:multiloopfin}
Sample topologies of diagrams that give UV finite contributions to the quasi-quark PDFs.}
\end{center}
\end{figure}
%%%%%%%%%%%%%%%%%%%%%%%%%%%%%%%

%%%%%%%%%%%%%%%%%%%%%%%%%%%%%%%
\begin{figure}[h]
\begin{center}
\includegraphics[width=3in]{sumSE}
\caption{\label{fig:sumSE}
Contributions to the general diagrams in Fig.~\ref{fig:multiloopdiv}(a) with all loop diagrams reorganized in terms of 1PI diagrams.}
\end{center}
\end{figure}
%%%%%%%%%%%%%%%%%%%%%%%%%%%%%%%
